\subsection{A fixed-point XGBoost reference}
\label{sec:fixedpoint}
As a first numerical implementation check, we emulate the XGBoost reference
with signed fixed-point inputs, thresholds and scores. A format with total
width $W_b$ and integer width $I_b$ has $F_b=W_b-I_b$ fractional bits and
representable range
$[-2^{I_b-1},\,2^{I_b-1}-2^{-F_b}]$. Define its quantizer by
\begin{equation}
 Q_{W_b,I_b}(u)=\operatorname{clip}_{W_b,I_b}
 \!\left(2^{-F_b}\operatorname{round}_{\rm even}(2^{F_b}u)\right),
 \label{eq:quantizer}
\end{equation}
where rounding chooses the nearest integer with ties to even, and clipping
saturates at the format's endpoints. We quantize inputs and thresholds before
evaluating the strict-left XGBoost cuts. Let $s_{t\ell}$ denote a scalar leaf contribution with the learning rate
absorbed, and $s_0$ the initial margin. These values
use the score format; tree scores are combined by a balanced adder tree with
casts to that format after additions. Its result is added to the quantized
initial margin and cast once more. No sigmoid is needed for selection.

The maximum absolute training input is 8.96 and the largest absolute tree
threshold is 4.73, so we use $I_b=5$ for inputs and thresholds. For the
$T=100$ tree scores, define
\begin{equation}
 A=|s_0|+\sum_{t=1}^{T}\max_\ell|s_{t\ell}|\simeq15.09.
 \label{eq:score-bound}
\end{equation}
At the coarsest fractional width $F_{\min}=4$, rounding the $T$ leaf values
and initial margin can add at most $(T+1)2^{-F_{\min}-1}$ to this bound.
Their quantized sum and all partial sums therefore fit within an absolute
bound of 18.25, covered by $I_b=6$ for scores. With no overflow, sums of values
on this common fixed-point grid introduce no further fractional rounding.
No validation input lies outside the chosen input range. We vary the common
fractional width while keeping these integer widths fixed. Table~\ref{tab:fixedpoint}
shows the resulting validation performance. The 17/18-bit input/score pair
retains the reference working-point efficiency; the coarsest 9/10-bit pair
loses performance and produces more substantial score ties. Small positive
changes in the interpolated point are not interpreted as an established
accuracy gain.

\begin{table}[htbp]\centering\small
\caption{XGBoost fixed-point emulation on the same validation partition.
Widths include the sign bit; input and threshold formats are identical.
The native reference has $\epsilon_s=55.48\%$ and AUC 0.92056.
The strict-threshold column uses background efficiency at most 1\%.}
\label{tab:fixedpoint}\setlength{\tabcolsep}{5pt}
\begin{tabular}{@{}rrrrr@{}}\toprule
Input bits & Score bits & $\epsilon_s$ [\%] & $\epsilon_s^{\rm strict}$ [\%] & AUC\\\midrule
23 & 24 & 55.48 & 55.48 & 0.92056 \\
17 & 18 & 55.49 & 55.49 & 0.92056 \\
15 & 16 & 55.46 & 55.45 & 0.92056 \\
13 & 14 & 55.56 & 55.53 & 0.92048 \\
11 & 12 & 55.47 & 55.29 & 0.92040 \\
9 & 10 & 52.05 & 50.85 & 0.91760 \\
\bottomrule
\end{tabular}
\end{table}
Independent C++ executions agree exactly for the tested fixed-point outputs.
This establishes numerical agreement of the emulators; FPGA resource use,
latency and initiation interval require synthesis for a specified device,
clock and interface.

\begin{implementation}
\textbf{Tested chain.} \file{docs/conifer_xgboost_v1/README.md} specifies the
EL9 container, CMSSW workspace, dependency pins and commands. Conversion and
HLS C++ generation use conifer 1.11; the axis-aligned reference is also evaluated
with the actual CMSSW conifer 1.7 external. Its bridge rejects oblique weights.
The full validation compares 13,280,000 native/exported leaf IDs, plus boundary
checks. Each of eight C++ profiles is checked on 512 HLS vectors, including all
324 frozen XGBoost handover vectors. Every fixed-point comparison is exact.
\file{fpga/conifer/compat.py} reconciles native binary32 thresholds, leaf
values, initial margin and actual depth; upstream conversion is retained
as a diagnostic. It is important to preserve XGBoost's input rounding and
strict-left equality convention. The score-type settings in the table are
\code{AP_RND_CONV, AP_SAT}; axis weights are exact one-hot constants.

\file{docs/conifer_xgboost_v1/validation.json} contains full provenance and
per-profile metrics. \file{scripts/validate_xgboost_conifer.py} reads training
inputs for range selection and validation arrays for evaluation; test arrays
are not loaded. Generated projects use conifer's default device/clock only as
export placeholders. No FPGA synthesis has been run, and CPU timings that
include file I/O are not reported as hardware latency.
\end{implementation}
